% POSC 315 - Syllabus - Online Asynchronous - Fall 2024

\documentclass[12pt, letterpaper]{article}
\usepackage[margin=1in]{geometry}
\usepackage{xcolor}
\usepackage[colorlinks=true, linkcolor=blue, urlcolor=blue, citecolor=blue]{hyperref}
\usepackage{url}
\usepackage{graphicx}
\usepackage{bookmark}
\usepackage{enumitem}
\frenchspacing
\usepackage{titlesec}
\titleformat{\section}{\normalfont\Large\bfseries}{\thesection.}{1em}{}
\usepackage[scaled]{helvet}
\renewcommand\familydefault{\sfdefault}

\begin{document}
\title{\includegraphics[width=6cm]{/home/snags/Nextcloud/Syllabi/Images/stacked.png} \\\vspace{1ex} \textbf{Introduction to Public Policy}}

\author{POSC 315-02 -- Fall 2024}
\date{Tuesdays and Thursdays at 11:30 in MH 554 \\\vspace{2ex} \textsc{August 26 -- December 13, 2024}}

    \maketitle


\subsection*{Professor: David P. Adams, Ph.D.}

\subsubsection*{Contact Information:}

\begin{itemize}
	\item Office: 516 Gordon Hall
	\item Phone/Text: \href{tel:+16572784770}{(657) 278-4770}
	\item Zoom: \href{https://fullerton.zoom.us/j/3347502369}{\texttt{https://fullerton.zoom.us/j/3347502369}}
	\item website: \href{https://dadams.io}{\texttt{dadams.io}}
	\item email: \href{dpadams@fullerton.edu}{\texttt{dpadams@fullerton.edu}}
	\item Office Hours:
        \begin{itemize}
            \item Tuesdays and Thursdays: 1:00pm -- 2:00pm; Thursdays: 5:00pm -- 6:00pm
            \item Schedule meetings throughout the week: \href{https://dadams.io/appt}{\texttt{dadams.io/appt}}
        \end{itemize}  
\end{itemize}


\section{Catalog Description}

	Federal domestic policymaking. Structure, functions, and relationships among American national institutions, including executive, legislative and judicial branches, media, political parties, and pressure groups.

\section{Course Description}

	In this course, students will explore and engage in thoughtful discussions on the processes and key players in creating public policy in the United States. The curriculum focuses on the structure, functions, and relationships among American national institutions, including the executive, legislative, and judicial branches of government, the media, political parties, and interest groups. We will examine the various official and unofficial influences on the policy process and the limitations imposed by institutional and structural factors.

	\noindent This course delves into the historical and constitutional development of the policy process, as well as the distinct characteristics of public policy within a federal system. By understanding policy-making in the context of a constitutional republic and a federalized governance system, students will gain a deeper appreciation for the complexities and uncertainties surrounding agenda setting, policy making, policy implementation, and policy evaluation in the American political system.

\section{Student Learning Objectives}

By the end of this course, students will be able to

\begin{enumerate}
	\item discuss and explain the key features of the public policy-making process in the United States;
	\item recognize and describe the distinct stages of the public policy process;
	\item describe the various internal and external actors that influence public policy, their interactions, and their impact on the policy process;
	\item articulate the historical and contemporary structures and institutions that facilitate, expand, or constrain the public policy process;
	\item differentiate and describe the various theories that attempt to explain the drivers and influences leading to public policy change or maintaining the status quo; and
	\item apply their knowledge of the policy process to analyze specific policy domains impacted by multiple policy actors and diverse elements of the policy process within the context of public policy-making in the United States.
\end{enumerate}


\section{General Education Information}

\subsection*{Requirements Satisfied}

	This course satisfies General Education Explorations in Social Sciences subarea D.5. The writing assignments in this course, including the policy memo papers and current event summaries described below, meet the requirement of UPS 411.201: 
	\begin{quote}Writing assignments in General Education courses shall involve the organization and expression of complex data or ideas and careful and timely evaluations of writing so that deficiencies are identified, and suggestions for improvement and/or for means of remediation are offered. Evaluations of the student's writing competence shall determine the final course grade\ldots .\end{quote}
    A grade of “D” (1.0) or higher is required to meet this General Education requirement. A grade of “D-“ (0.7) or below will not satisfy this General Education requirement.

\subsection*{General Education Student Learning Goals}

	Students completing courses in this subarea shall encounter the following learning goals:

\begin{enumerate}
	\item Examine problems, issues, and themes in the social sciences in greater depth; in a variety of cultural, historical, and geographical contexts; and from different disciplinary and interdisciplinary perspectives.
	\item Analyze and critically evaluate the application of social science concepts and theories to particular historical, contemporary, and future problems or themes, such as economic and environmental sustainability, globalization, poverty, and social justice.
	\item Analyze and critically evaluate constructs of cultural differentiation, including ethnicity, gender, race, class, and sexual orientation, and their effects on the individual and society.
	\item Apply theories and concepts from the social sciences to address historical, contemporary, and future problems confronting communities at different geographical scales, from local to global.
\end{enumerate}

\section{Text Books}

\begin{itemize}
\item Kraft, Michael E. and Scott R. Furlong. 2025. \emph{Public Policy: Politics, Analysis, and Alternatives}. 8th ed. CQ Press: Thousand Oaks, Calif.

\item Smith, Kevin B. and Christopher W. Larimer. 2017. \emph{The Public Policy Theory Primer}. 3rd ed. Westview Press: Boulder, Colo.
\end{itemize}
\section{Technical Competencies}

Students are expected to have the following technical competencies to succeed in this course:
\begin{itemize}
    \item Basic computer skills, including the ability to navigate the internet, use email, and create and save documents.
    \item Proficiency in using \emph{Canvas}, including submitting assignments, participating in discussions, and accessing course materials.
    \item Ability to use word processing software, such as Microsoft Word or Google Docs, to create and format documents.
    \item Access to a reliable computer and internet connection.
    \item Ability to use Zoom for virtual meetings and discussions.
    \item Basic knowledge of online communication tools, such as email and discussion boards.
    \item Ability to use online research tools and databases to find academic sources.
\end{itemize}

\section{Technical Problems}

\subsection*{University IT Help Desk}

Contact the instructor immediately to document the problem if you encounter any technical difficulties. Then contact the \href{http://www.fullerton.edu/it/students/helpdesk/index.php}{Student IT Help Desk} for assistance. You can also call the Student IT Help Desk at \href{tel:+16572788888}{(657) 278-8888}, \href{mailto:StudentITHelpDesk@fullerton.edu}{email}, visit them at the Pollak Library North \href{http://www.fullerton.edu/it/students/sgc/index.php}{Student Genius Center}, or log on to the \href{http://my.fullerton.edu/}{my.fullerton.edu} portal and click ``Online IT Help'' followed by ``Live Chat''.

\subsection*{Canvas Support}

If you encounter any technical difficulties with Canvas, call the Canvas Support Hotline at \href{tel:+18553027528}{855-302-7528}, visit the \href{https://community.canvaslms.com/docs/DOC-10720-67952720329}{Canvas Community}, or click the ``Help'' button in the lower left corner of Canvas and select ``Report a Problem''. The \href{https://cases.canvaslms.com/liveagentchat?chattype=student&sfid=001A000000YzcwQIAR}{Student Support Live Chat} is available 24 hours a day, 7 days a week.

\section{Response Time} I will strive to respond to all student emails, Discord posts, and \emph{Canvas} messages within 24 hours, except on weekends and holidays. If you are still awaiting a response within 24 hours, please send a follow-up message. If you are still waiting to receive a response within 48 hours, please send another follow-up message and contact me via phone or text at \href{tel:+16572784770}{(657) 278-4770}.


\section{University Student Policies}

In accordance with UPS 300.00, students must be familiar with certain policies applicable to all courses. Please review these policies as needed and visit this Cal State Fullerton website \texttt{\href{https://fdc.fullerton.edu/teaching/student-info-syllabi.html}{https://fdc.fullerton.edu/teaching/student-info-syllabi.html}} for links to the following information:

\begin{enumerate}
    \item   University learning goals and program learning outcomes.
    \item	Learning objectives for each General Education (GE) category.
    \item	Guidelines for appropriate online behavior (netiquette).
    \item	Students' rights to accommodations for documented special needs.
    \item   Campus student support measures, including Counseling \& Psychological Services, Title IV and Gender Equity, Diversity Initiatives and Resource Centers, and Basic Needs Services.
    \item   Disability Support Services (DSS) information.
    \item	Academic integrity (refer to UPS 300.021).
    \item	Actions to take during an emergency.
    \item	Library services information.
    \item	Student Information Technology Services, including details on technical competencies and resources required for all students.
    \item	Software privacy and accessibility statements.
\end{enumerate}

\section{Course Student Policies}

\subsection*{Course Communication}
All course announcements and communications will be sent via \emph{Canvas} and university email. Students are responsible for regularly checking their \emph{Canvas} notifications and email. Students are also responsible for ensuring that their \emph{Canvas} notifications are set to receive messages from the course. Students are expected to check \emph{Canvas} and their email at least once daily.

\subsection*{Due Dates}
All assignments are due by 11:59 p.m. on the specified due date. Save for extenuating circumstances, late assignments will not be accepted unless prior arrangements have been made with the professor. Students are responsible for submitting assignments on time. Failure to submit an assignment on time may result in a failing grade for the assignment. 

\subsection*{Alternative Procedures for Submitting Work}
Students are expected to submit all assignments via \emph{Canvas}. If you cannot submit an assignment via \emph{Canvas}, please get in touch with the professor to discuss alternative submission procedures.

\subsection*{Retention of Student Work}
Students are responsible for retaining copies of all assignments submitted for this course. Students are also responsible for retaining copies of all graded assignments returned by the professor.

\subsection*{Extra Credit}
 There is one bonus assignment available in this course. A documentary analysis assignment will be available in Week 15. The assignment is optional and will provide extra credit points on the final exam. The assignment will be graded based on the quality of the analysis and the depth of engagement with the course material.

\subsection*{Academic Integrity}
Students are expected to adhere to the highest standards of academic integrity. Any student found to have engaged in academic dishonesty will be subject to the sanctions described in the \href{https://www.fullerton.edu/senate/publications_policies_resolutions/ups/UPS%20300/UPS%20300.021.pdf}{Academic Dishonesty Policy} (UPS 300.021). Academic dishonesty includes, but is not limited to, cheating, plagiarism, fabrication, facilitating academic dishonesty, and submitting previously graded work without prior authorization. Students are expected to be familiar with the university's policy on academic dishonesty and to adhere to this policy in all aspects of this course. Any student who has questions about the policy should ask the professor for clarification.

\subsection*{Written Work}
All written work must be submitted in a professional format, including proper grammar, spelling, and punctuation. Written work must also be properly cited using the appropriate citation style. Students are expected to follow the guidelines for written work provided by the professor and to seek clarification if they have questions about the requirements.

\subsection*{Plagiarism}
Plagiarism is a serious violation of academic integrity and will not be tolerated in this course. Plagiarism includes, but is not limited to, copying and pasting text from sources without proper citation, paraphrasing text from sources without proper citation, and submitting work that is not your own. Students are expected to properly cite all sources used in their work and to submit original work. Failure to do so may result in a failing grade for the assignment and further disciplinary action. Written work will be checked for plagiarism using plagiarism detection software.

\subsection*{AI Generated Text} 
Large language models, such as OpenAI's ChatGPT-4o, Anthropic's Claude Opus 3, and others, have made it easier to generate text that mimics human writing. While these models can be useful for generating ideas and content, they can also be misused to produce work that is not original. Students are expected to use AI-generated text responsibly and to ensure that all work submitted in this course is their own. Failure to do so may result in a failing grade for the assignment and further disciplinary action. Written work will be checked for the use of AI-generated text using plagiarism detection software.

\subsection*{Participation}

Students are expected to participate in all course activities. This includes completing all assigned readings, watching all assigned videos, and participating in all discussions. Students are expected to participate in discussions in a professional and respectful manner. Students are expected to be familiar with the university policy on netiquette and to adhere to this policy in all aspects of this course. Any student who has questions about the policy should ask the professor for clarification. 

\subsection*{Netiquette}
Students are expected to adhere to the university's policy on netiquette. The university's policy on netiquette is as follows:
\begin{quote}Netiquette refers to a set of behaviors that are appropriate for online activity (e.g., social media, email, discussions, presentations). All personnel at Cal State Fullerton are expected to demonstrate appropriate online behavior at all times. A good summary of netiquette can be found in the \href{https://canvashelp.fullerton.edu/m/Student/l/1336786-student-what-is-netiquette}{CSUF Canvas self-help guides}, which adapts ten rules to the online course situation from the website for the book \href{http://www.albion.com/netiquette/corerules.html}{Netiquette by Virginia Shea} and other sources referenced at the bottom of the guide.\end{quote}


\section{Course Assignment Descriptions}
\subsection*{Discussion Posts (20\% of Final Grade)} 
\begin{itemize}
    \item \textbf{Description}: \emph{Students will participate in \underline{six of the eight} online discussion boards}, responding to prompts related to the current week's readings and lectures. Prompts may involve analyzing case studies, applying theoretical concepts to current events, or debating policy issues.
    \item \textbf{Expectations}: Students are expected to post a response to the prompt and engage with their classmates' posts. Discussion posts are due by 11:59 p.m. on Saturdays for the week they are assigned. Late posts will not be accepted.
    \item \textbf{Objectives}: Engage with course material, apply theoretical concepts to real-world scenarios, and develop critical thinking skills.
    \item \textbf{Grading Criteria}: Posts will be graded based on the quality of the response, depth of analysis, timeliness of contributions, and engagement with classmates' posts.
\end{itemize}

\subsection*{Policy Analysis Paper (35\% of Final Grade)}

This course includes a scaffolded and peer-reviewed writing project where students will conduct a comprehensive analysis of a policy issue. The project will be completed in stages, with opportunities for peer review and feedback at each stage.

\subsubsection*{Project Stages}
\begin{enumerate}
    \item \textbf{Topic Selection and Proposal (Week 3)}
    \begin{itemize}
        \item Submit a one-page proposal outlining the policy issue, its significance, and a preliminary research question.
        \item Participate in peer reviews of proposals.
    \end{itemize}
    \item \textbf{Annotated Bibliography (Week 6)}
    \begin{itemize}
        \item Compile an annotated bibliography of 8-10 sources, summarizing and evaluating each source's relevance to your research question.
        \item Exchange and review peers' bibliographies.
    \end{itemize}
    \item \textbf{Literature Review (Week 9)}
    \begin{itemize}
        \item Write a 3-4 page literature review synthesizing key findings and identifying gaps in the research.
        \item Review and comment on peers' literature reviews.
    \end{itemize}
    \item \textbf{Draft Analysis (Week 12)}
    \begin{itemize}
        \item Submit a 6-8 page draft analysis, including an introduction, literature review, policy analysis, and preliminary conclusions.
        \item Conduct detailed peer reviews using a provided rubric.
    \end{itemize}
    \item \textbf{Final Paper (Week 15)}
    \begin{itemize}
        \item Submit a final 10-12 page policy analysis paper, incorporating peer feedback and additional research or revisions.
    \end{itemize}
\end{enumerate}

This project will help you develop your research, analytical, and writing skills while providing a deep understanding of policy theory and the policy process.

\subsection*{Midterm Exam (20\% of Final Grade)}
An in-class midterm exam will be held in Week 8. The exam will cover material from the first half of the course, including readings, lectures, and discussions. The exam will consist of multiple-choice, short answer, and essay questions.

\subsection*{Final Exam (25\% of Final Grade)}
An in-class final exam will be held during the scheduled final exam period. The final exam will be comprehensive and cover all material from the course, including readings, lectures, discussions, and assignments. The final exam will consist of multiple-choice, short answer, and essay questions.

\section{Grading}
Grades will be allocated based on the following components:
\begin{itemize}
    \item \textbf{Discussion Posts}: 20\%
    \item \textbf{Policy Analysis Paper}: 35\%
    \item \textbf{Midterm Exam}: 20\%
    \item \textbf{Final Exam}: 25\%    
\end{itemize}

\subsection*{Grading Scale}

\begin{itemize}
    \item A+: 97.00 -- 100.00
    \item A: 93.00 -- 96.99
    \item A-: 90.00 -- 92.99
    \item B+: 87.00 -- 89.99
    \item B: 83.00 -- 86.99
    \item B-: 80.00 -- 82.99
    \item C+: 77.00 -- 79.99
    \item C: 73.00 -- 76.99
    \item C-: 70.00 -- 72.99
    \item D+: 67.00 -- 69.99
    \item D: 63.00 -- 66.99
    \item D-: 60.00 -- 62.99
    \item F: 0.00 -- 59.99
\end{itemize}

\section{Course Schedule}

\begin{itemize}
    \item \textbf{Week 1: August 27 and August 29}
    \begin{itemize}
        \item \textbf{Topic:} Introduction to Public Policy
        \item \textbf{Readings:} Smith \& Larimer, Chapter 1; Kraft \& Furlong, Chapter 1
        \item \textbf{Discussion:} Why Study Public Policy?
    \end{itemize}

    \item \textbf{Week 2: September 3 and September 5}
    \begin{itemize}
        \item \textbf{Topic:} Government Institutions and Policy Actors
        \item \textbf{Readings:} Kraft \& Furlong, Chapter 2; Smith \& Larimer, Chapter 3
        \item \textbf{Discussion:} Decision-Making Models and Institutional Rational Choice
    \end{itemize}

    \item \textbf{Week 3: September 10 and September 12}
    \begin{itemize}
        \item \textbf{Topic:} Understanding Public Policymaking
        \item \textbf{Readings:} Kraft \& Furlong, Chapter 3; Smith \& Larimer, Chapter 5
        \item \textbf{Assignment:} Topic Selection and Proposal Due
    \end{itemize}

    \item \textbf{Week 4: September 17 and September 19}
    \begin{itemize}
        \item \textbf{Topic:} Policy Design and Typologies
        \item \textbf{Readings:} Smith \& Larimer, Chapter 4; Kraft \& Furlong, Chapter 3
        \item \textbf{Discussion:} Policy Typologies as Analytic Tools
    \end{itemize}

    \item \textbf{Week 5: September 24 and September 26}
    \begin{itemize}
        \item \textbf{Topic:} Policy Analysis Introduction
        \item \textbf{Readings:} Kraft \& Furlong, Chapter 4; Smith \& Larimer, Chapter 6
        \item \textbf{Discussion:} Rationalist and Post-Positivist Approaches
    \end{itemize}

    \item \textbf{Week 6: October 1 and October 3}
    \begin{itemize}
        \item \textbf{Topic:} Impact Analysis and Program Evaluation
        \item \textbf{Readings:} Smith \& Larimer, Chapter 7; Kraft \& Furlong, Chapter 6
        \item \textbf{Assignment Due:} Annotated Bibliography 
    \end{itemize}

    \item \textbf{Week 7: October 8 and October 10}
    \begin{itemize}
        \item \textbf{Topic:} Policy Implementation
        \item \textbf{Readings:} Smith \& Larimer, Chapter 8; Kraft \& Furlong, Chapter 6
        \item \textbf{Discussion:} Implementation Analysis and Challenges
    \end{itemize}

    \item \textbf{Week 8: October 15 and October 17}
    \begin{itemize}
        \item \textbf{Topic:} Midterm Review
        \item \textbf{Exam:} Midterm Exam
    \end{itemize}

    \item \textbf{Week 9: October 22 and October 24}
    \begin{itemize}
        \item \textbf{Topic:} Economic and Budgetary Policy
        \item \textbf{Readings:} Kraft \& Furlong, Chapter 7
        \item \textbf{Assignment:} Literature Review Due
    \end{itemize}

    \item \textbf{Week 10: October 29 and October 31}
    \begin{itemize}
        \item \textbf{Topic:} Social and Environmental Policy
        \item \textbf{Readings:} Kraft \& Furlong, Chapters 8 and 11
        \item \textbf{Assignment:} Policy Issue Project Proposal due
    \end{itemize}

    \item \textbf{Week 11: November 5 and November 7}
    \begin{itemize}
        \item \textbf{Topic:} Policy Change and Innovation
        \item \textbf{Readings:} Smith \& Larimer, Chapter 9
        \item \textbf{Discussion:} Policy Change and Innovation
    \end{itemize}

    \item \textbf{Week 12: November 12 and November 14}
    \begin{itemize}
        \item \textbf{Topic:} Foreign Policy and Homeland Security
        \item \textbf{Readings:} Kraft \& Furlong, Chapter 12
        \item \textbf{Assignment:} Draft Analysis Due
    \end{itemize}

    \item \textbf{Week 13: November 19 and November 21}
    \begin{itemize}
        \item \textbf{Topic:} Policy Conflicts and Strategies
        \item \textbf{Readings:} Kraft \& Furlong, Chapter 13
        \item \textbf{Discussion:} Policy Conflicts and Strategies
    \end{itemize}

    \item \textbf{Break}: November 26 and November 28
    \begin{itemize}
        \item \textbf{Note:} Thanksgiving Break (No Class on November 26 and 28)
    \end{itemize}

    \item \textbf{Week 14: December 3 and December 5}
    \begin{itemize}
        \item \textbf{Topic:} Class Wrap-Up and Bonus Documentary Screening 
        \item \textbf{Discussion:} Course Reflections and Final Thoughts
    \end{itemize}

    \item \textbf{Week 15: December 10 and December 12}
    \begin{itemize}
        \item \textbf{Topic:} Final Exam Review
        \item \textbf{Assignment:} Final Paper Due
        \item \textbf{Bonus Assignment:} Optional Documentary Analysis Due
    \end{itemize}

    \item \textbf{Week 16: December 17}
    \begin{itemize}
        \item \textbf{Final Exam: 11:00 - 12:50}
    \end{itemize}
\end{itemize}


\end{document}
